\section{Рабочий проект}

\subsection{Спецификация программных компонентов}

В данном разделе приведено детальное описание интерфейсов разработанных функций и структур данных.

\textbf{Функция расширения ключа (key-expansion)}
\begin{itemize}
	\item входные параметры: key — исходный ключ шифрования (вектор 16 байт);
	\item выходные данные: вектор раундовых ключей;
	\item назначение: генерация полного набора ключей из одного секретного ключа, необходимых для всех раундов шифрования.
\end{itemize}

\textbf{Функция дополнения данных (pkcs7-pad)}
\begin{itemize}
	\item входные параметры: data — исходный вектор байтов произвольной длины;
	\item выходные данные: новый вектор, длина которого кратна 16 байтам;
	\item назначение: подготовка данных к блочному шифрованию путем добавления байтов-заполнителей согласно стандарту PKCS\#7.
\end{itemize}

\textbf{Функция удаления дополнения (pkcs7-unpad)}
\begin{itemize}
	\item входные параметры: data — вектор расшифрованных байтов;
	\item выходные данные: вектор байтов исходной длины;
	\item назначение: удаление байтов-заполнителей после дешифрования для восстановления оригинального сообщения.
\end{itemize}

\textbf{Функция шифрования блока (encrypt-block)}
\begin{itemize}
	\item входные параметры: data — блок данных (16 байт), key — ключ;
	\item выходные данные: зашифрованный блок (16 байт);
	\item назначение: реализация 10 раундов AES (SubBytes, ShiftRows, MixColumns, AddRoundKey).
\end{itemize}

\textbf{Функция шифрования данных (encrypt-data)}
\begin{itemize}
	\item входные параметры: data — вектор исходных данных произвольной длины, key — ключ;
	\item выходные данные: вектор зашифрованных данных;
	\item назначение: основной цикл шифрования, выполняет дополнение данных (pad), разбивку на блоки и применение функции encrypt-block.
\end{itemize}

\textbf{Функция дешифрования блока (decrypt-block)}
\begin{itemize}
	\item входные параметры: data — зашифрованный блок (16 байт), key — ключ;
	\item выходные данные: расшифрованный блок (16 байт);
	\item назначение: реализация обратных преобразований AES для одного блока.
\end{itemize}

\textbf{Функция дешифрования данных (decrypt-data)}
\begin{itemize}
	\item входные параметры: data — вектор зашифрованных данных, key — ключ;
	\item выходные данные: вектор исходных данных;
	\item назначение: основной цикл дешифрования, выполняет обработку блоков функцией decrypt-block и финишное удаление дополнения (unpad).
\end{itemize}

\textbf{Функция запуска тестов (run-tests)}
\begin{itemize}
	\item входные параметры: отсутствуют (используются встроенные тестовые сообщения);
	\item выходные данные: вывод результатов в консоль (сравнение оригинальных и дешифрованных данных);
	\item назначение: проверка корректности всей библиотеки путем шифрования и дешифрования сообщений разной длины.
\end{itemize}

\subsection{Модульное тестирование функций библиотеки}

Для проверки корректности работы алгоритма был разработан модуль тестирования run-tests. На рисунке \ref{code:tests} представлен фрагмент кода сценария тестирования (проверка длинного массива данных).

\begin{figure}[H]
	\begin{lstlisting}[language=Lisp]
		(defun run-tests ()
		(let ((key #(1 2 3 4 5 6 7 8 9 10 11 12 13 14 15 16)))
		
		;; ... (Тесты 1 и 2 опущены) ...
		
		;; Тест 3: данные больше размера блока
		(let* ((data #(10 20 30 40 50 60 70 80 90 100
		110 120 130 140 150 160 170 180 190))
		(enc (encrypt-data data key))
		(dec (decrypt-data enc key)))
		(print "========================")
		(print "Test 3: long data")
		(print "Original:") (print data)
		(print "Encrypted:") (print enc)
		(print "Decrypted:") (print dec)
		(print "OK?")
		(print (equalp data dec)))))
	\end{lstlisting}  
	\caption{Фрагмент кода модуля тестирования}
	\label{code:tests}
\end{figure}

\subsection{Системное тестирование разработанной библиотеки}

Тестирование проводилось в интерактивной среде Common Lisp (SBCL). Были проверены три сценария работы программы, покрывающие различные варианты длины входных данных.

\textbf{Тест №1. Данные меньше размера блока.}
На рисунке \ref{test1:image} показан результат обработки массива длиной 2 байта. Программа корректно дополнила данные до 16 байт (PKCS\#7), зашифровала и восстановила оригинал.

\begin{figure}[H]
	\center{\includegraphics[width=0.9\linewidth]{test_short}}
	\caption{Результат теста №1: Данные меньше размера блока}
	\label{test1:image}
\end{figure}

\textbf{Тест №2. Данные, равные размеру блока.}
На рисунке \ref{test2:image} показан результат для данных длиной ровно 16 байт. Согласно стандарту, был добавлен новый блок заполнения.

\begin{figure}[H]
	\center{\includegraphics[width=0.9\linewidth]{test_block}}
	\caption{Результат теста №2: Данные равные размеру блока}
	\label{test2:image}
\end{figure}

\textbf{Тест №3. Данные больше размера блока.}
На рисунке \ref{test3:image} показана обработка массива, превышающего размер блока (19 байт). Алгоритм корректно обработал данные в цикле, расшифровка совпала с оригиналом.

\begin{figure}[H]
	\center{\includegraphics[width=0.9\linewidth]{test_long}}
	\caption{Результат теста №3: Данные больше размера блока}
	\label{test3:image}
\end{figure}